% !TEX root = ../main.tex

\chapter{系统设计与实验分析}

\section{系统总体架构}

为了验证本文所提方法的可行性与有效性，本章设计并实现了一套面向金融大数据的智能风控与决策分析系统。系统整体架构自底向上划分为数据层、算法层、服务层与应用层四个层次：

\begin{enumerate}
    \item \textbf{数据接入与治理层}：负责异构金融数据（交易日志、征信记录、网络行为特征）的高吞吐接入、特征工程提取与隐私脱敏。
    \item \textbf{核心算法层}：集成图神经网络（GNN）、集成树模型（XGBoost/LightGBM）与大语言模型检索增强（RAG）管道，实现欺诈团伙识别与高危交易判定。
    \item \textbf{微服务调度层}：基于高性能 RPC\footnote{远程过程调用（Remote Procedure Call, RPC）允许程序通过网络调用其他进程或主机提供的服务。}框架实现毫秒级推理调用与高并发流量治理。
    \item \textbf{业务可视化层}：为业务风控专家提供可视化审查工作台与全链路溯源图谱。
\end{enumerate}

关于残差网络、注意力机制与自适应优化的工作，可联合引用\cite{ResNet,Transformer,Adam}，三个连续编号将自动合并为范围。在模型实现中，深层神经网络的结构设计可参考残差连接与批归一化方法\cite{ResNet,BatchNormalization}，参数优化可参考 Adam 算法\cite{Adam}。这些引用展示了英文文献及多篇文献在正文中的标注方式。

\section{实验评估与结果对比}

\subsection{数据集与评价指标}
实验选取公开的金融信用卡反欺诈基准数据集（Credit Card Fraud Detection）以及某股份制商业银行真实脱敏反洗钱数据集进行交叉验证。评估指标涵盖精确率（Precision）、召回率（Recall）、受试者工作特征曲线下面积（ROC-AUC）以及综合度量指标（F1-Score）。

\subsection{消融实验分析}
通过逐步剥离不同核心特征模块，考察各模块对整体性能提升的边际贡献。实验结果表明，融合多跳图拓扑结构特征与上下文语义特征后，系统在假阳性率（FPR）控制在 0.1\% 的严苛条件下，召回率相较传统基线方法提升了 12.4\%。

这充分证明了深度图结构表征与大模型语义增强在复杂金融风险识别中的优越性。
